\documentclass[11pt,a4paper]{article}

% -------------------------------------------------------------------------
%  PACKAGES
% -------------------------------------------------------------------------
\usepackage[T1]{fontenc}
\usepackage[utf8]{inputenc}
\usepackage{mathpazo}
\usepackage{amsmath,amssymb}
\usepackage{geometry}
\usepackage{xcolor}
\usepackage{titlesec}
\usepackage{enumitem}
\usepackage{booktabs}
\usepackage{array}
\usepackage{multirow}
\usepackage{tabularx}
\usepackage{tcolorbox}
\tcbuselibrary{skins,breakable}
\usepackage{fancyhdr}
\usepackage{hyperref}
\usepackage{caption}
\usepackage{tikz}
\usetikzlibrary{arrows.meta,positioning}
\usepackage[numbers,sort&compress]{natbib} % or \usepackage{cite}
\bibliographystyle{unsrt}                   % or plain / apalike

\geometry{a4paper,left=2.6cm,right=2.6cm,top=2.6cm,bottom=2.8cm}

% -------------------------------------------------------------------------
%  PALETTE  (kept consistent with the companion slide deck)
% -------------------------------------------------------------------------
\definecolor{Ink}{HTML}{203247}
\definecolor{Accent}{HTML}{2E6E8E}
\definecolor{Highlight}{HTML}{C5681B}
\definecolor{AccentPale}{HTML}{EAF2F5}
\definecolor{Paper}{HTML}{FBFBF9}
\definecolor{Muted}{HTML}{5F6F7B}
\definecolor{Rule}{HTML}{C9D6DC}

\pagecolor{Paper}
\color{Ink}

\hypersetup{
  colorlinks=true,
  linkcolor=Accent,
  citecolor=Accent,
  urlcolor=Accent,
  pdftitle={Rooting for Kepler: A Comparative Benchmark of Newton-Type Solvers},
  pdfauthor={CSE 402 Subsection A2}
}

% -------------------------------------------------------------------------
%  SECTION TITLING — small-caps, ruled, numbered in a muted serif
% -------------------------------------------------------------------------
\titleformat{\section}
  {\normalfont\Large\scshape\color{Ink}}
  {\textcolor{Accent}{\thesection}}{0.6em}{}
  [\vspace{-0.3em}{\color{Rule}\titlerule[0.8pt]}]
\titlespacing*{\section}{0pt}{1.6em}{0.9em}

\titleformat{\subsection}
  {\normalfont\large\scshape\color{Accent}}
  {\thesubsection}{0.55em}{}
\titlespacing*{\subsection}{0pt}{1.1em}{0.4em}

% -------------------------------------------------------------------------
%  HEADER / FOOTER
% -------------------------------------------------------------------------
\pagestyle{fancy}
\fancyhf{}
\renewcommand{\headrulewidth}{0.5pt}
\renewcommand{\footrulewidth}{0pt}
\fancyhead[L]{\footnotesize\scshape\color{Muted}CSE 402 \textbar\ Numerical Methods}
\fancyhead[R]{\footnotesize\scshape\color{Muted}Project Proposal}
\fancyfoot[C]{\footnotesize\color{Muted}\thepage}
\renewcommand{\headrule}{\hbox to\headwidth{\color{Rule}\leaders\hrule height \headrulewidth\hfill}}

% -------------------------------------------------------------------------
%  SMALL UTILITIES
% -------------------------------------------------------------------------
\newcommand{\key}[1]{\textcolor{Highlight}{\bfseries #1}}
\newcommand{\marginlab}[1]{\textcolor{Accent}{\scshape\footnotesize #1}}

\newtcolorbox{note}{
  enhanced, breakable, colback=AccentPale, colframe=Rule,
  boxrule=0.5pt, arc=0pt, outer arc=0pt,
  left=8pt, right=8pt, top=6pt, bottom=6pt,
  before skip=1em, after skip=1em
}
\newtcolorbox{notetitled}[1]{
  enhanced, breakable, colback=AccentPale, colframe=Rule,
  boxrule=0.5pt, arc=0pt, outer arc=0pt,
  left=8pt, right=8pt, top=6pt, bottom=6pt,
  before skip=1em, after skip=1em,
  fonttitle=\scshape\footnotesize\color{Accent},
  colbacktitle=AccentPale, coltitle=Accent,
  title={#1}
}

\newcommand{\dividerline}{\vspace{0.2em}{\color{Rule}\hrule height 0.5pt}\vspace{0.6em}}

\renewcommand{\arraystretch}{1.3}

\setlist[itemize]{label=\textcolor{Accent}{\small$\bullet$}, leftmargin=1.4em, itemsep=0.15em}
\setlist[enumerate]{leftmargin=1.6em, itemsep=0.2em}

% =========================================================================
\begin{document}

% -------------------------------------------------------------------------
%  TITLE BLOCK
% -------------------------------------------------------------------------
\begin{titlepage}
\pagecolor{Paper}
\thispagestyle{empty}
\begin{center}
  \vspace*{1.6cm}
  {\color{Accent}\scshape\large CSE 402 \,\textbar\, Numerical Methods\, --\, Project Proposal}\\[0.6cm]
  {\color{Rule}\rule{0.55\textwidth}{0.8pt}}\\[1.1cm]

  {\color{Ink}\scshape\Huge Rooting for Kepler}\\[0.35cm]
  {\color{Muted}\Large\itshape a comparative benchmark of\\[2pt] Newton-type root finders and starting guesses}\\[1.4cm]

  {\color{Rule}\rule{0.32\textwidth}{0.6pt}}\\[1.1cm]

  \begin{minipage}{0.72\textwidth}
    \centering
    \small\color{Muted}
    An empirical study of classical, domain-specialised, and\\
    2024--2025 with-memory Newton-type variants applied to\\
    Kepler's equation, evaluated on RadVel's radial-velocity fitting pipeline.
  \end{minipage}

  \vspace{2.2cm}
  {\scshape\color{Accent}\footnotesize Subsection A2 \quad\textbar\quad Student IDs}\\[0.35cm]
  {\color{Ink}2105036 \quad 2105050 \quad 2105053 \quad 2105056 \quad 2105059}

  \vfill
  {\footnotesize\color{Muted}\today}
\end{center}
\end{titlepage}

\setcounter{page}{1}

% =========================================================================
\section{Motivation}

Radial-velocity orbit fitting reduces, at its numerical core, to a single
repeated task: solving the transcendental Kepler equation
\[
  f(E) = E - e\sin E - M = 0
\]
for the eccentric anomaly $E$, given a mean anomaly $M$ and eccentricity
$e$. A fitting pipeline such as \textsc{RadVel} calls this solver on the
order of $10^5$--$10^7$ times over the course of a single MCMC run, so the
choice of root finder is not a cosmetic detail --- it is one of the
dominant costs in the fit, and its accuracy floor propagates directly into
the recovered orbital parameters.

\textsc{RadVel} itself uses \key{Danby's} higher-order iterative method
rather than an ordinary Newton--Raphson update, which already signals that
the "obvious" solver is not automatically the best one for this equation.
At the same time, the general numerical-analysis literature on
Newton-type iteration has continued to move: 2024--2025 work on
\key{with-memory} multipoint schemes reports polynomial convergence orders
in excess of $10$, obtained by recycling already-computed function values
across iterations rather than by paying for new derivative evaluations.

\begin{notetitled}{Why this is not a settled question}
None of the recent with-memory papers were written with Kepler's equation
in mind, and Kepler's equation has a cost structure --- one transcendental
call can cheaply return \emph{both} $\sin E$ and $\cos E$ --- that generic
efficiency-index comparisons ($E=\rho^{1/\gamma}$) do not capture. Whether
a higher formal convergence order actually pays off once this
Kepler-specific evaluation cost is accounted for, and across the
$(e,M)$ operating range that real radial-velocity fits encounter, is an
open, testable question rather than a foregone conclusion.
\end{notetitled}

% =========================================================================
\section{Scope Statement}

The purpose of this project is to build a controlled, Kepler-specific
benchmark of five root-finding variants and one alternative starting-guess
strategy, and to measure --- not assume --- how they compare.

\begin{tcolorbox}[
  enhanced, colback=Paper, colframe=Accent, boxrule=0.9pt, arc=0pt,
  left=10pt,right=10pt,top=8pt,bottom=8pt, before skip=1em, after skip=1.2em
]
\scshape\color{Ink}\small We are explicitly \emph{not} setting out to prove
that the newer with-memory Newton-type variants outperform, or
underperform, the classical solvers already used for Kepler's equation.
\normalfont\upshape\color{Ink} The project is designed as an experiment
whose outcome is not fixed in advance: the newer variants may win on some
axes (order, iteration count) and lose on others (per-iteration cost,
robustness at high eccentricity), and part of the deliverable is
characterising exactly where that trade-off falls.
\end{tcolorbox}

% =========================================================================
\phantomsection
\section*{Addressing Supervisor Feedback}
\begingroup\emergencystretch=1.6em

An earlier draft of this proposal was scoped narrowly, around substituting an
ordinary Newton--Raphson update for the Danby iteration currently used in
\textsc{RadVel}. Supervisory review identified that this framing was too
restrictive. The standard Newton--Raphson method is unlikely to be the right
replacement for a domain-tuned higher-order scheme, and a study built around a
single alternative could not establish whether any observed difference was
attributable to the iteration itself. The revisions below were made in
response, and they account for the present structure of Sections~3 and~4.

\vspace{0.5em}
\begingroup
\small
\setlength{\tabcolsep}{5.5pt}
\renewcommand{\arraystretch}{1.15}
\begin{tabularx}{\textwidth}{@{}>{\raggedright\arraybackslash}X >{\raggedright\arraybackslash}X >{\raggedright\arraybackslash}p{1.3cm}@{}}
\toprule
\textbf{Feedback} & \textbf{Revision} & \textbf{Section} \\
\midrule
Do not limit the study to standard Newton--Raphson.
& Candidate set widened to five solvers across three categories.
& \textcolor{Accent}{3} \\
\addlinespace[0.55em]
Explore a range of alternatives, not one substitution.
& Two with-memory schemes (2024--2025) added.
& \textcolor{Accent}{3.2} \\
\addlinespace[0.55em]
Retain the baseline so results stay interpretable.
& Newton--Raphson kept as anchor; category sizes balanced.
& \textcolor{Accent}{3.1} \\
\addlinespace[0.55em]
Do not presuppose that the newer methods are superior.
& Scope statement fixes the design, not the expected result.
& \textcolor{Accent}{2, 6} \\
\addlinespace[0.35em]
\bottomrule
\end{tabularx}
\endgroup

\vspace{0.4em}
\begin{note}
The net effect of these revisions is a change of research question. The study
no longer asks whether one specific solver can replace another. It asks under which
conditions of eccentricity, mean anomaly, and Kepler-specific evaluation cost
each class of Newton-type method is preferable. This is a question whose
answer is not determined in advance by the choice of methods.
\end{note}
\endgroup

% =========================================================================
\section{Candidate Methods}

Five solvers are selected to span three categories that are each
represented by exactly two members, plus one initial-guess strategy that
is evaluated orthogonally to all of them.

\subsection{Selection rationale}
\begin{itemize}
  \item \key{Baseline:} ordinary Newton--Raphson, to anchor every
        comparison against the textbook method.
  \item \key{Established Kepler-specific solvers:} two variants that are
        already used or benchmarked against in the Kepler-equation
        literature, representing the current practical standard rather
        than a strawman.
  \item \key{Recent with-memory variants (2024--2025):} two multipoint
        Newton-type schemes that raise convergence order via
        Hermite-interpolated self-accelerating parameters, at no
        additional evaluation points per iteration.
  \item \key{Initial-guess layer:} a symbolic-regression-derived starting
        formula, applied in front of every iterative solver above,
        isolated as its own experimental factor.
\end{itemize}

\subsection{Method table}
\begingroup
\small
\begin{tabularx}{\textwidth}{@{}>{\raggedright\arraybackslash}p{2.05cm} >{\raggedright\arraybackslash}p{3.55cm} >{\raggedright\arraybackslash}p{1.1cm} >{\raggedright\arraybackslash}X@{}}
\toprule
\textbf{Category} & \textbf{Method} & \textbf{Order} & \textbf{Role in the benchmark} \\
\midrule
Baseline & Newton--Raphson & 2 & Textbook anchor; one evaluation point, full $\sin,\cos$ pair per iteration. \\[0.5em]
\multirow{2}{2.05cm}{\raggedright Established Kepler solvers}
 & Danby's method & 4 & RadVel's production solver; already domain-tuned to this equation. \\
 & Markley's method & \textemdash & Closed-form / near-analytic quartic solver widely cited as a Kepler-specific standard. \\[0.5em]
\multirow{2}{2.05cm}{\raggedright With-memory (2024--2025)}
 & NWM9 (Mittal, Panday \& J\"antschi, 2024) & 8.8989 & Single self-accelerating parameter on an optimal 8th-order base; controlled, lower-order sibling of NWM11. \\
 & NWM11 (Mittal, Panday, J\"antschi \& Bolundu\c{t}, 2025) & 10.7446 & Bi-parametric variant; primary "recent" candidate --- 3 evaluation points, only 1 full $\sin,\cos$ pair per iteration. \\[0.5em]
Guess layer & Napier's guess (2024, arXiv:2411.15374) & n/a & Symbolic-regression starting formula; drop-in replacement for the simple/canonical/RadVel starts, used in front of every solver above. \\
\bottomrule
\end{tabularx}
\endgroup

\vspace{0.4em}
\begin{note}
Two established, two new, and one baseline gives every iterative solver a
same-size comparison group, while the guess-layer factor is deliberately
kept independent so that "better start" and "better iteration" are never
conflated in the results.
\end{note}

% =========================================================================
\section{Evaluation Framework}

\subsection{Kepler-specific cost accounting}
Raw convergence order is not, by itself, a fair currency for comparing
these methods on this particular equation. A single \texttt{sincos} call
returns $\sin E$ and $\cos E$ together at roughly the cost of one
transcendental evaluation, so a scheme that needs both quantities at one
point is markedly cheaper than one that needs $\sin$ and $\cos$ at two
different points. Every method above will therefore be costed by:
\begin{enumerate}
  \item the number of \emph{distinct} evaluation points used per
        iteration;
  \item how many of those points require a full $\sin,\cos$ pair versus
        $\sin$ (or $\cos$) alone; and
  \item whether any derivative at a new point is synthesised from
        divided-difference / Hermite machinery instead of a fresh
        transcendental call.
\end{enumerate}

\subsection{Benchmark grid}
Each solver--guess combination will be run over a broad $(e,M)$ grid,
including the ordinary operating region of typical radial-velocity fits
and the classical pathological corner --- high eccentricity, $M$ near
zero --- where Newton-type iteration for Kepler's equation is known to be
most fragile.

\subsection{Metrics}
\begin{itemize}
  \item \key{Empirical convergence order}, measured directly from the
        residual sequence rather than taken on faith from the source
        paper.
  \item \key{Iterations to tolerance} and \key{wall-clock time} per
        solve, and per full RadVel fit.
  \item \key{Residual and reference-solution error} against a
        high-precision root.
  \item \key{Robustness}: failure or non-convergence rate across the
        $(e,M)$ grid.
  \item \key{Downstream effect}: propagation of solver error through
        $M \rightarrow E \rightarrow \nu \rightarrow v_r$ into the fitted
        parameters $P$, $e$, and $K$, and comparison against RadVel's MCMC
        posterior scale.
\end{itemize}

% =========================================================================
\section{Work Plan}

\begin{enumerate}
  \item \textbf{Implementation.} Re-derive and implement all five solvers
        and the Napier guess in a common scalar-arithmetic framework,
        instrumented to record evaluation-point-level cost, not just
        iteration count.
  \item \textbf{Verification.} Validate each implementation against its
        source paper's reported convergence order on a small controlled
        test set before running the full grid.
  \item \textbf{Grid benchmark.} Run every solver--guess combination
        across the $(e,M)$ grid; record all metrics in
        Section~4.3.
  \item \textbf{Error propagation study.} Vary solver tolerance and trace
        the resulting error through to refit orbital parameters on a real
        radial-velocity dataset.
  \item \textbf{Monte Carlo uncertainty comparison.} Generate noisy
        realisations of the same dataset, refit, and compare
        measurement-noise-driven parameter spread against solver-induced
        shifts and RadVel's MCMC posterior.
  \item \textbf{Synthesis.} Report where each method category wins,
        loses, or ties, once Kepler-specific evaluation cost is
        accounted for.
\end{enumerate}

% =========================================================================
\section{Expected Outcomes}

\begin{itemize}
  \item A reusable, cost-instrumented benchmark harness comparing
        classical, domain-specialised, and 2024--2025 with-memory
        Newton-type solvers on Kepler's equation specifically.
  \item An empirical answer --- not an assumed one --- to whether the
        higher formal order of recent with-memory schemes survives
        contact with Kepler's specific $\sin,\cos$ cost structure and
        RadVel's real operating range of $(e,M)$.
  \item A quantified picture of how solver-level numerical error
        propagates into fitted orbital parameters, set against RadVel's
        own MCMC uncertainty for context.
\end{itemize}

\vspace{1.4em}
\dividerline

\section*{References}
\begingroup
\small
\begin{enumerate}[label={[\arabic*]}, leftmargin=2.2em, itemsep=0.7em]

  \item B. J. Fulton, E. A. Petigura, S. Blunt, and E. Sinukoff,
  ``RadVel: The Radial Velocity Fitting Toolkit,''
  \textit{Publications of the Astronomical Society of the Pacific},
  vol. 130, no. 986, p. 044504, 2018.
  doi: \href{https://doi.org/10.1088/1538-3873/aaaaa8}
  {10.1088/1538-3873/aaaaa8}.

  \item J. M. A. Danby,
  ``The solution of Kepler's equation. I,''
  \textit{Celestial Mechanics},
  vol. 40, nos. 3--4, pp. 303--312, 1987.
  doi: \href{https://doi.org/10.1007/BF01230252}
  {10.1007/BF01230252}.

  \item F. L. Markley,
  ``Kepler's Equation Solver,''
  \textit{Celestial Mechanics and Dynamical Astronomy},
  vol. 63, no. 1, pp. 101--111, 1995.
  doi: \href{https://doi.org/10.1007/BF00691917}
  {10.1007/BF00691917}.

  \item R. Mittal, A. Panday, and L. J{\"a}ntschi,
  ``An Optimal Higher-Order Derivative-Free with-Memory Iterative Scheme
  for Solving Nonlinear Equations,''
  \textit{Symmetry},
  vol. 16, no. 5, p. 588, 2024.
  doi: \href{https://doi.org/10.3390/sym16050588}
  {10.3390/sym16050588}.

  \item R. Mittal, A. Panday, L. J{\"a}ntschi, and L.-M. Bolundu{\c{t}},
  ``High-Efficiency With-Memory Iterative Solvers for Nonlinear Equations:
  Development, Theoretical Analysis, and Numerical Validation,''
  \textit{Mathematics},
  vol. 13, no. 2, p. 245, 2025.
  doi: \href{https://doi.org/10.3390/math13020245}
  {10.3390/math13020245}.

  \item K. Napier,
  ``Analytic Starting Guesses for Kepler's Equation via Symbolic Regression,''
  \textit{arXiv preprint arXiv:2411.15374}, 2024.
  doi: \href{https://doi.org/10.48550/arXiv.2411.15374}
  {10.48550/arXiv.2411.15374}.

\end{enumerate}
\endgroup
```

\end{document}